\section{Related Work}
\label{sec:related}

\subsection{Visual representations for robotic manipulation}
Vision foundation models have become widely used for robotic representation learning. Self-supervised approaches extract robust features from static observations via time-contrastive learning (e.g., R3M~\cite{nair2022r3m}, VIP~\cite{ma2022vip}, LIV~\cite{ma2023liv}) and masked autoencoders (e.g., MAE~\cite{he2022masked}, MVP~\cite{radosavovic2023real}, VC-1~\cite{majumdar2023we}). To incorporate semantic priors, vision-language models align visual features with language instructions, enabling open-vocabulary manipulation (e.g., CLIPort~\cite{shridhar2022cliport}, RT-1~\cite{brohan2022rt}). 


These methods produce static representations that capture \textit{what} is in the scene, but not \textit{how} it changes over time. To extract more robust and task-relevant signals, recent approaches extend these static representations using the Information Bottleneck (IB) principle~\cite{tishby2000information}. By strictly limiting information capacity, IB-based methods learn compressed representations that filter out irrelevant distractors, improving downstream reinforcement learning~\cite{you2024multimodal} and behavior cloning~\cite{bai2025rethinking}. However, these IB-based methods still lack explicit modeling of task-relevant dynamics, as they compress the current state while not predicting its future evolution.

\subsection{Predictive modeling for visual control}
An alternative approach is learning representations by predicting future observations, which effectively captures environment dynamics. Several approaches focus on pixel-level video prediction (e.g., RoboNet~\cite{dasari2019robonet}, UniPi~\cite{du2023learning}, SUSIE~\cite{black2023zero}), for visual planning or sub-goal generation. Other approaches, such as Voltron~\cite{karamcheti2023language} and LaVA-Man~\cite{zhu2025lava}, extend such pixel-level prediction by integrating language and action conditioning into image reconstruction and forward video prediction. However, forecasting at the pixel level is computationally expensive and overly sensitive to task-irrelevant visual distractions such as lighting variations or background movement~\cite{lecun2022path,bardes2024revisiting}.

Joint-Embedding Predictive Architectures (JEPAs)~\cite{lecun2022path} address the limitation of pixel-level prediction by shifting the prediction objective to the latent space, which improves efficiency. I-JEPA~\cite{assran2023self} and V-JEPA~\cite{bardes2024revisiting} learn powerful spatial and spatiotemporal representations by predicting the latent embeddings of masked regions. In the robotics domain, MC-JEPA~\cite{benchetrit2024mcjepa} jointly models motion and content features, which can capture both scene appearance and dynamics. As a recent extension to predictive models, ToBo (Token Bottleneck)~\cite{kim2025token} squeezes the current scene into a compact bottleneck token to predict future latent patches, forcing the model to encode temporal dependencies. 

Despite these advances, standard latent predictive models (e.g., Siamese MAE~\cite{gupta2023siamese}) still attempt to forecast the \textit{full} future state, inevitably entangling controllable transitions with uncontrollable passive dynamics. In contrast, our proposed TRAIL combines the efficiency of latent prediction with a task-relevant bottleneck. Unlike prior methods, we explicitly predict the \textit{residual} between current and future states, providing policies with a gradient-like directional cue that isolates task-relevant dynamics.

\CO{A related way to exploit action-free video is to \emph{invent} an action space. Latent action
models infer a latent action $u_t$ from an observation pair with an inverse dynamics model, ask a
forward model to reconstruct the future observation from $u_t$, and force $u_t$ through a
low-capacity bottleneck so that it retains the controllable part of the
transition~\cite{schmidt2024learning, bruce2024genie, ye2025latent}. A policy is pre-trained to emit
$u_t$ on unlabelled video, and a small action-labelled set then aligns the latent actions with real
commands; the same recipe has been scaled to large robot corpora~\cite{agibot2025world}. TRAIL
shares the premise that action-free video contains usable dynamics, but not the mechanism. We place
the capacity constraint on \emph{visual memory access} rather than on an action variable, and the
quantity we extract, $\Delta_z$, lives in the observation space of a frozen perceptual encoder, so
no latent-to-real action alignment stage is required. This also avoids an embodiment mismatch: a
latent action inferred from human hands is defined relative to the human embodiment, whereas a
state residual is a direction in a perceptual space shared by human and robot. The trade-off is
that $\Delta_z$ is not an action and therefore cannot be rolled out for planning, which latent
action models can do; the two lines are thus complementary, ours acting as a representation-side
prior rather than a substitute for the action space.}
